\chapter{Space Age Texas}

The Cold War and Space Race were an environment where industrial and military imperatives, academic breakthroughs, and shrewd political maneuvering converged, creating the conditions for a world-class computational facility to take root in Austin. At the same time, the demands of the oil industry were driving advanced computing nearby in Houston. In fact, the technological stories in the neighboring cities were in many ways intertwined, and it makes sense to view them in a unified way, especially as the national political role of Texas grew over the same years. A surprising aspect of all this is how little popular impact it had at the time, perhaps because it was simply overshadowed by the traditional, larger than life, Texas mythology. Equally, Texas society and culture have always been firmly rooted in the land itself, and in practical concerns. Texas computing developed as a means to an end in very targeted and specific applications essential for working the land or for defending it. Despite this popular amnesia, much can be said for preserving the history of advanced computing in Austin and Houston from the fifties up through the seventies, the era of the mainframes, focusing on the practical nature of the work that was done in numerical modeling and simulation. It’s enlightening to start nearby with the human computers at Los Alamos and their rapid adoption of the forerunner of the IBM CPC, follow that thread onward to both Santa Monica and Houston, then onward to Austin, looking at a handful of real breakthroughs along the way: Simplex, ADI, SOR, subsurface, and orbit determination. It’s worth noting that these developments were not foreseen and not popularly recognized at the time or since, not just in Texas but generally. Science fiction and popular science promoted space travel and artificial brains, but not the somewhat invisible and mathematical world of modeling and simulation. The role of Texas in these developments is well worth clarifying and preserving.

From a national and political standpoint, the era of Lyndon Johnson’s influence was a driving force behind federal investment in Texas, as he was instrumental in drafting and enacting the National Aeronautics and Space Act of 1958, then chairing the National Aeronautics and Space Council. Johnson strategically recognized that the space race was a unique opportunity to modernize the economic infrastructure of Texas, and his ties to the oil, aerospace, and electronics industries helped steer contracts and development to the region. The establishment of the Manned Spacecraft Center in Houston was a signature achievement of this effort, anchoring the Apollo program in Texas. The congressional committee in charge of NASA’s funding was chaired by Houston’s representative. Combined with a 1,000-acre land donation from Rice University, nearby air operations at Ellington Air Force Base, and access to Galveston Bay for transporting heavy spacecraft components, Houston was politically and logistically ideal. This brought millions in initial federal construction contracts to the region, creating a permanent demand for advanced engineering and computational capabilities. 
Undoubtedly UT Austin would have played an important regional role in any case, but Johnson’s influence on the course of events can not be overlooked. As one example, the UT Astronomy department was able to build a major new telescope at McDonald Observatory in West Texas during the Johnson years, the 107-inch telescope completed in 1968, ostensibly in support of the NASA Apollo program. It’s questionable whether that would have happened without his active involvement. A similar effect was seen in the engineering departments on campus in Austin, especially those directly involved in the Space Race and connected with NASA and the missile programs. In fact, somewhat of a campus revolution took place in engineering and computing under the impulse of federal funding.

At the same time the new telescope was under construction in West Texas, a specially designed home for the new CDC 6600 supercomputer was under construction beside the UT Tower and Main Building, at the head of the East Mall. This unique site still bears the simple name of Computation Center, though it has long since become obsolete and been repurposed as a backup emergency data center. It is built of the local limestone, like most of the classic campus buildings, but is mostly underground, with the Main Building’s wide eastern terrace as its roof. Some say it was designed not to block views of the Main Building and Tower. Others say that the motivation was the CDC 6600’s cooling requirements and its sheer complexity. The system contained 400,000 individual transistor components and more than 100 miles of internal wiring. Cooling was via a Freon refrigerant system that circulated through metal plates in contact with tightly packed circuit boards designed to keep wire lengths short and signal speeds high.

From the Computation Center’s terrace roof, a wide flight of stairs descends to the East Mall and directly faces a unique landmark complex of buildings in the distance, the Lyndon B. Johnson Presidential Library. It’s an ideal vista, as the Library is a superb embodiment of the fifties and sixties era and a monument to the spirit of the times. Essentially, it’s a living Star Trek set and one can easily imagine Captain Kirk and Mr. Spock beaming down beside the Library’s main Tower, a fittingly Space Age monument. 

Long before President Kennedy famously assembled his inner circle of Ivy League social scientists, physical scientists in the United States were already being treated as "the best and the brightest" by military leaders and heads of state who actively sought their counsel. This elite status was the driving force behind the rapid academic expansion at institutions like UT Austin. Because the military and political establishment believed these scientists were essential for securing national defense and winning the Space Race, massive federal budgets were redirected to cultivate this talent pool. 
At UT Austin, this resulted in significantly increased defense funding of higher education, especially through the establishment of facilities like the Defense Research Laboratory and eventually the Center for Space Research, which concentrated elite talent to solve complex orbital, aerospace, and computing challenges. As they were integrated into the national security apparatus to execute Big Science projects, the life of the isolated university scholar was replaced by what was described as a bureaucratic rat race. These top-tier academics increasingly spent their time managing federal defense contracts, directing massive teams of assistants, and sitting on government advisory committees. This all reflected a profound faith that elite scientific minds, armed with state-of-the-art supercomputers and federal patronage, could engineer solutions to the nation's most critical geopolitical and technical problems. 

To align with this new national mandate, the university officially renamed its aeronautical engineering program to the Department of Aerospace Engineering in 1959, the exact same year NASA commenced its full operations. By 1968, it had merged with Engineering Mechanics to form the modern, comprehensive department. The university became a center of excellence in gravity field modeling and orbit determination, skills fundamental not only to the space and missile programs but also to regional economic drivers like Texas oil exploration. 

The career of professor Byron Tapley is an excellent example of the real-world impact of this period. He earned his bachelors in mechanical engineering at UT in 1956, just before the launch of Sputnik, followed by his doctorate in engineering mechanics in 1960, and he joined the UT faculty in 1959 and quickly became a pillar of the university's aerospace ambitions. In 1961, he established UT's orbit mechanics program, and by 1963, he was teaching the first orbital mechanics courses on campus, focusing heavily on low-thrust trajectory analysis and optimization. He subsequently steered the department's growth as its chair from 1966 to 1977. He was one of the first professors in the new Aerospace Engineering and Engineering Mechanics department that formed with lightning speed in its own new building on what was then the north-west corner of campus. 

Tapley went on to lead the Center for Space Research, working in a host of overlapping fields including gravity mapping, geology, oceanography, oil exploration, navigation, and astronomy, and was active well into the next century. In 2014, it was possible to interact on a daily basis with a leading member of the first generation of the Space Age, a role-model at the frontier of engineering for well over fifty years who had effectively begun his career with Sputnik in 1957. Tapley’s fifty-plus-year career produced a staggering human legacy that rippled throughout the national economy. He supervised hundreds of students who went on to form the backbone of the American aerospace and engineering sectors. [[sat2]]
Throughout the late fifties and sixties, the Los Angeles Model of aerospace development, characterized by a dense proximity of airframe manufacturers and federal research labs, began to reach geographic and economic saturation. This catalyzed a dispersal of elite researchers and Big Science methodologies from Southern California defense industries into emerging Sunbelt academic institutions like UT Austin. The concentration of defense research at Los Alamos and White Sands also played a role in this story. By importing technology and scientists directly from the defense world, Texas broke the geographic determinism that had previously clustered advanced scientific work almost exclusively on the coasts.The strategic convergence of elite talent, state-of-the-art supercomputing hardware, and targeted federal funding allowed UT to institutionalize its expertise. 

The creation of the UT Computation Center established UT as a national center for computing, aerospace research, and numerical analysis. From this base, it projected its influence directly into the nation's most critical Cold War aerospace and defense research programs. Because these programs focused heavily on the complex trajectory analysis and gravity field modeling essential for space missions, UT Austin forged a powerful connection with federal laboratories like NASA's Jet Propulsion Laboratory in Los Angeles. By the seventies, there were well known concentrations of UT alumni in the national labs, particularly Los Alamos, NASA JPL, and NASA Goddard Space Flight Center. At JPL, the so-called Texas Mafia of UT alumni became legendary. The same can be said for the major Exxon offices in Houston, where UT alumni helped create significant aspects of the Exxon corporate culture. 

In essence, from the sixties onward UT Austin played a major role on the national computing, engineering, and scientific stage, much as universities such as MIT and Stanford already had done beginning with WWII and the Cold War. Ultimately, UT integrated its academic mission into the national security apparatus. By blending the abstract rigors of the academy with the industrial demands of the military, UT capitalized on targeted federal funding to help create a Texas technopolis, an ecosystem where the boundaries between military necessity, academic research, and industrial profit were intentionally and permanently blurred.

\section{The Defense Education Act And NASA}

The 1958 National Defense Education Act and National Aeronautics and Space Act operated in tandem as the legislative catalysts for the new era, permanently altering the trajectory of American universities. Following the 1957 launch of Sputnik, the U.S. recognized an existential need to rapidly mobilize its scientific and industrial resources, transforming higher education into an explicit instrument of national policy. This era formalized the Space Race and established the federal government's significant role in education through student loans, fellowships, and curriculum support. It marked the first time the federal government injected significant funding into higher education, framing it as a matter of national security. 

The results were transformative, creating the modern American university system and the federal student loan infrastructure. It effectively militarized the justification for education funding, leading to a golden age of American research universities. It converted the U.S. education system from a local responsibility into a key component of national security policy. And the resonance with the simultaneous apogee in Washington DC of a native Central Texan with very strong ties to Austin in some sense created a perfect storm on the UT campus. Several different forces were coming together to form a period at UT that launched its engineering programs from the era of desktop mechanical calculators to the forefront of computational science and supercomputing within a decade. 

It is true that defence electronics companies such as Texas Instruments and aerospace giants like Convair were nearby in Dallas, but few would have predicted in 1958 that Austin and Houston would become centers of computing power and home of the manned spaceflight program. The success of Tracor, the city's first homegrown Fortune 500 company, demonstrated the viability of commercializing academic research in acoustics and defense electronics. This momentum was institutionalized by the arrival of IBM and Texas Instruments, which provided the industrial scale necessary to support a permanent high-tech ecosystem. Unlike the centralized federal model seen in Houston with NASA, the Austin tech scene of the sixties was a product of academic cooperation, setting the stage for the city's later evolution as a primary center for semiconductor design and software innovation. 

At the same time, driven by the G.I. Bill and intense federal defense spending, graduate enrollments in physics skyrocketed. Between 1945 and 1950, the rate of new physics doctorates grew at nearly twice the rate of all other fields combined, eventually tripling prewar highs by the Korean War and tripling again after Sputnik. To meet this challenge, the National Science Foundation poured unprecedented funding into the development of new undergraduate and precollege curricula across physics, mathematics, chemistry, and engineering. The most notable example was the Physical Science Study Committee, spearheaded by MIT physicist Jerrold Zacharias. A deepened understanding of reality is in fact one of the types of knowledge that was clarified and transferred by the new academic channels created during this period. The structure and content of the textbooks, courses, and even departments were reshaped in the fifties and sixties. The fact they have been fairly stable for the following decades points out how definitive the changes were.

The National Science Foundation poured unprecedented resources into the development of new instructional materials. These reforms aimed to dismantle the traditional cookbook laboratory experience, replacing it with a pedagogy that emphasized the process of scientific inquiry and the underlying structural logic of physical laws. The most enduring monument to this pedagogical shift is found in The Feynman Lectures on Physics, which reconceptualized the introductory physics sequence for the modern era. By integrating quantum mechanics and electromagnetism into the foundational curriculum, Richard Feynman and his colleagues at Caltech challenged the classical hierarchy of scientific instruction. Feynman’s approach mirrored the broader Cold War trend of democratizing high-level theoretical concepts for a burgeoning workforce of aerospace engineers and research physicists. 

The Space Race was an expression of a much broader cultural and ideological zeitgeist that defined the fifties and sixties. This era was defined by a profound faith in big science, systems analysis, and centralized, computer-driven problem-solving. It was also marked by an equally dramatic transformation of this worldview in the subsequent decade, with the seventies becoming an era of limits and inward-looking Earth-centered humanism. The meaning of Big Science was tangible to everyone experiencing the birth of nuclear weaponry via the Manhattan Project, the resulting national labs at Los Alamos and Livermore, the intercontinental missile and manned spaceflight programs, the SAGE continental air defence system, jet aircraft, gargantuan particle accelerators such as SLAC and Fermi Lab, globe circling nuclear submarines, and the so-called electronic brains of IBM marketing campaigns. Through these massive initiatives, research methodologies shifted from isolated academic efforts to hierarchical, government-funded teamwork, deeply integrating universities like UT Austin into the national security state. 

Ironically, the pinnacle of Big Science, the Apollo space program, helped catalyze the inward-looking environmentalism of the seventies. The famous 1968 Apollo 8 Earthrise photograph taken by astronaut William Anders provided humanity with a profound new perspective of Earth as a fragile, finite inheritance that needed to be handled with utmost care. The backlash was just as dramatic as Sputnik, with Earth days, silent springs, population bombs, limits to growth, and an overall surge of ecological and humanistic pursuits. 

These phases were readily apparent in Austin and Houston, and in a particularly concentrated form on the UT campus and within the UTCC, where the swing from crew-cuts and ties to long hair, beards, and sandals took place with dizzying speed. The UT administration physically insulated its most valuable, commercially-oriented programs. The engineering complex was deliberately built on the north-east outskirts of the campus, isolating it from the central liberal arts buildings with a built environment that actively discouraged student protest. Ultimately, the campus rebellions of this era set the stage for the entrepreneurialization of the university, a strategic pivot where institutions embraced overt corporate partnerships, commercialized research, and market-driven metrics to regain financial control and bypass the spaces carved out by students during the seventies. 

\section{Big Science And Engineering}

This was also a period of geographic and political realignment in the American scientific-industrial complex. There was a purposeful migration of leading researchers from defense-oriented corporations and laboratories back into universities. This institutionalized knowledge transfer was particularly evident in the evolution of aerospace and electrical engineering, where specialists who had actively been working in defence related projects were expected to return to purely civil academic careers and spread the technologies that they had been exposed to in the defense world, and this career pattern is clearly evident among prominent engineers and scientists of the era. Faculty members arriving from defense contractors or national laboratories like Los Alamos brought with them an unprecedented level of federal funding and a mission-driven research culture. 

This shift transformed the university ecosystem from traditional isolated scholarship to Big Science, where academic success was increasingly measured by the ability to manage large-scale, multi-disciplinary projects. This was the conscious means by which sensitive technologies were released into the broader commercial and educational spheres, via the individuals directly involved, making the interpersonal links between leading researchers central to understanding what was taking place. Communications were still relatively slow, with a reliance on physical mail, telephone, and travel via railway and later commercial aviation. Electronic and digital communications were relatively rare or not yet existent. 
As the traditional clustering of scientific activity on the Northeast and California coasts, especially with regard to the Los Angeles Model of aerospace development, reached geographic and economic saturation, Texas capitalized on the federal mandate to decentralize national security infrastructure. The state escaped the era's geographic determinism by intentionally importing technology and elite scientists directly from the defense establishment. There’s no more clear sign of this than the fact that advanced technological work gained a foothold and then took on a life of its own in Texas, perhaps one of the harshest environments for pursuits not based on land ownership and natural resources. 

A prime example of this institutionalized knowledge transfer was the migration of researchers from the Harvard Underwater Sound Laboratory to UT Austin. Following World War II, researchers like Paul Boner utilized their defense experience to establish the Defense Research Laboratory, later the Applied Research Laboratory, at UT Austin to perform advanced military technology research for the Navy. These migrating specialists brought with them the Big Science paradigm, which permanently shifted the academic model from isolated, individual scholarship to one characterized by hierarchical organization, massive research teams, and a heavy reliance on federal funding. Consequently, the traditional academic lifestyle was increasingly replaced by a bureaucratic rat race of managing government contracts, directing crews of technicians, and serving on advisory committees. Because electronic and digital communications were still rudimentary, the transfer of sensitive defense technologies into the educational sphere required the physical relocation and interpersonal collaboration of the individuals involved. 

This created a direct pipeline of ideas and techniques between established defense hubs and emerging Texas institutions. For instance, a DRL team led by George Wood collaborated with the Navy Electronics Laboratory in San Diego to improve minehunting sonar, subsequently bringing advanced methods for recording and analyzing data back to the Texas research environment, where acoustics and seismic data were already critical for subsurface exploration. Ties between California and Texas were reinforced in aerospace engineering where the growth of what would become the Center for Space Research, led by Byron Tapley, forged such a powerful interpersonal and institutional link between UT Austin and the NASA Jet Propulsion Laboratory in Los Angeles that the resulting pipeline of talent became legendary within the aerospace industry. 

The emergence of the American technopolis in the mid-twentieth century was not a product of isolated regional growth but rather the result of highly integrated, multi-nodal networks of experts, methodologies, and capital flowing between primary strategic hubs. During the fifties and sixties, the connection between Austin and Los Angeles served as a vital axis for the development of the nation's aerospace, electronic, and strategic capabilities. The population of Los Angeles County surged from 2.7 million in 1940 to 4.1 million in 1950, a 49\% increase driven by the explosive growth of the aerospace and defense industries. By the sixties, the region was the nation's primary manufacturing hub for high-tech weaponry and spacecraft. Austin, while smaller in scale, followed a similar trajectory. 

The transformation of the city from a regional agricultural and political center was made possible by the technology-transfer projects between the university and the defense industry. The presence of high-tech employers like Tracor and ARL:UT, which provided work experience to more than 4,000 students over its history, created a self-sustaining cycle of innovation and talent retention. Big science was taking on bureaucratic form in a proliferation of research centers and, at UT Austin, physical form in the establishment of the Pickle Research Campus, isolating sensitive Cold War projects far from the main campus. J.J. Pickle was at the heart of Lyndon Johnson’s inner-circle, and the name is at least a fitting reminder of the era. 

The political and geographic transformations driven by the Johnson era set the stage for a series of highly consequential technological decisions that would define the university's future role in national research. The creation of the UTCC and its associated academic centers was part of the national effort to ensure that the practical knowledge forged by Cold War and Space Race efforts was communicated into new academic disciplines, creating a lasting foundation for future generations of scientists and engineers.

This period saw UT Austin transform from a regional educational center into a critical component of the national defense infrastructure, largely through its interactions with Los Angeles-based entities such as the RAND Corporation, Ramo-Wooldridge (later TRW), and the United States Air Force Ballistic Missile Division. The architecture of the American Cold War was fundamentally dependent on the symbiotic relationship between the RAND Corporation and TRW, two organizations that transformed military procurement into a rigorous scientific discipline. RAND and TRW were the intellectual and industrial engines of the American Cold War. While the military provided the funding and the mandate, these two organizations provided the strategy and the systems engineering that built the U.S. nuclear arsenal and space program. 
Together, these entities created a closed loop of strategic planning and technical execution that allowed the United States to rapidly scale its nuclear and aerospace capabilities during the post-Sputnik era. Their early histories are deeply intertwined with the Air Force’s desire to harness scientific brainpower outside of the traditional military hierarchy. Safe in their Santa Monica headquarters, RAND civilians pioneered Operations Research, Game Theory, and Systems Analysis. If RAND provided the theory, TRW provided the management for Los Angeles based contractors such as Lockheed, North American Aviation, Douglas Aircraft, Northrop, Hughes Aircraft to build the hardware, creating the modern discipline of Systems Engineering. This institutionalized think tank model successfully merged the abstract rigors of the academy with the industrial demands of the military, cementing a technocratic governance structure that continues to dominate the American scientific and defense infrastructure. 

Two researchers, George Dantzig at RAND and David Young at TRW, demonstrate how these currents flowed into and shaped the ensuing growth of Texas and UT Austin. The raw power of early digital computers could only be unlocked through the creation of systematic, repeatable algorithms and a transition from theoretical exploration to large-scale numerical modeling. This shift was catalyzed by the pragmatic requirements of researchers like Dantzig and Young, who were tasked with managing unprecedentedly complex systems of equations. At the RAND Corporation, George Dantzig spearheaded the development of linear programming and the Simplex method, which were becoming essential for optimizing complex logistical and strategic military operations. The monumental task of actually implementing Dantzig's abstract mathematical algorithms on early machines fell to pioneering programmers such as William Orchard-Hays and Leola Cutler. Simultaneously, David Young was advancing iterative methods at TRW and later at UT Austin.

Dantzig’s development of the Simplex algorithm and Young’s work on iterative methods marked a fundamental change in approach for problem-solving, with a decisive move away from an adhoc, human-centric art, towards the standardized, machine-executable automated science of rigorous mathematical procedures. The intuitive calculation methods of the pre-computer era had to be replaced with standard algorithms that could be expressed in computer code and then executed automatically, a shift essential for solving problems at a scale previously unimaginable. 

The stories of both Dantzig and Young were also intertwined with the evolution of floating-point hardware and computerized linear algebra. Floating-point hardware developed fitfully across the fifties and sixties. It was not at all obvious, at least in the beginning, that floating-point would play the critical role that it did. In fact, prominent theorists famously opposed putting effort into floating-point hardware. That story will be explored here as an example of how the course of events was not at all obvious and straightforward to the researchers involved at the time, and to highlight the pragmatic roles played by Dantzig and Young. This period saw the birth of modern linear algebra, numerical analysis, scientific computing, and computational science. In the fifties and sixties, this field was largely synonymous with the complex challenge of solving systems of linear equations, using applied linear algebra executed with pure floating-point arithmetic on the most advanced machines available. Later, the overarching terms supercomputing and high-performance computing would also become synonymous with this field.
